\section{Revisión de contenidos}

\begin{enumerate}
  \item ¿Cuáles son los cuatro componentes principales de un computador y cuál es la función de cada uno?
  \item ¿Qué diferencia existe entre los registros visibles para el usuario y los registros de control y estado?
  \item ¿Qué función cumplen el Program Counter (PC), el Instruction Register (IR) y la Program Status Word (PSW) durante la ejecución de un programa?
  \item ¿En qué consiste el ciclo de instrucción y cuáles son sus etapas principales?
  \item ¿Qué tipos de acciones puede especificar una instrucción de máquina?
  \item ¿Qué es una interrupción y por qué mejora el aprovechamiento del procesador?
  \item ¿Cuáles son las cuatro clases principales de interrupciones? Mencione un ejemplo de cada una.
  \item Describa brevemente cómo se procesa una interrupción desde que ocurre hasta que el programa reanuda su ejecución.
  \item ¿Qué estrategias existen para gestionar múltiples interrupciones y cuáles son sus ventajas y desventajas?
  \item ¿Por qué los computadores utilizan una jerarquía de memoria en lugar de un único tipo de memoria?
  \item ¿Cómo está organizada una jerarquía de memoria típica, desde el nivel más rápido hasta el más lento?
  \item ¿Qué función cumple la memoria caché dentro de la jerarquía de memoria?
  \item Compare las técnicas de E/S programada, E/S mediante interrupciones y acceso directo a memoria (DMA).
  \item ¿Por qué DMA resulta más eficiente que las otras técnicas de entrada/salida para transferencias de grandes bloques de datos?
\end{enumerate}

\subsection*{Respuestas}

\paragraph{Respuesta 1} Procesador (ejecuta instrucciones), memoria principal (almacena programas y datos), módulos de entrada/salida (comunican el computador con el exterior) e interconexión del sistema (buses que permiten la comunicación entre los componentes).

\paragraph{Respuesta 2} Los registros visibles para el usuario almacenan datos y direcciones utilizados por los programas. Los registros de control y estado contienen información necesaria para coordinar la ejecución del procesador y el sistema operativo.

\paragraph{Respuesta 3}
\begin{itemize}
  \item PC: dirección de la siguiente instrucción.
  \item IR: instrucción actualmente ejecutada.
  \item PSW: estado del procesador (flags, modo de ejecución, interrupciones, etc.).
\end{itemize}

\paragraph{Respuesta 4} El ciclo de instrucción consta de la captación (fetch) de la instrucción, su ejecución (execute) y la comprobación de interrupciones antes de continuar con la siguiente instrucción.

\paragraph{Respuesta 5}
\begin{itemize}
  \item Transferencia procesador-memoria.
  \item Transferencia procesador-E/S.
  \item Procesamiento de datos.
  \item Control del flujo del programa.
\end{itemize}

\paragraph{Respuesta 6} Una interrupción es un evento que obliga al procesador a suspender temporalmente el programa en ejecución para atender otro evento. Permite que el procesador continúe trabajando mientras los dispositivos de E/S realizan operaciones lentas.

\paragraph{Respuesta 7}
\begin{itemize}
  \item Programa.
  \item Temporizador.
  \item Entrada/Salida (I/O).
  \item Fallo de hardware.
\end{itemize}

\paragraph{Respuesta 8} El procesador termina la instrucción actual, guarda el contexto del programa, ejecuta la rutina de servicio (ISR), restaura el contexto y continúa la ejecución desde el punto donde fue interrumpido.

\paragraph{Respuesta 9}
\begin{itemize}
  \item Deshabilitar interrupciones mientras se ejecuta una ISR (simple pero sin prioridades).
  \item Asignar prioridades, permitiendo interrupciones anidadas (nested interrupts), donde una ISR de mayor prioridad puede interrumpir otra de menor prioridad.
\end{itemize}

\paragraph{Respuesta 10} Porque no existe una memoria que sea simultáneamente muy rápida, de gran capacidad y de bajo costo. La jerarquía busca equilibrar estas tres características.

\paragraph{Respuesta 11} Registros \(\to\) Caché \(\to\) Memoria principal \(\to\) Almacenamiento secundario (SSD/HDD).

\paragraph{Respuesta 12} Almacena temporalmente los datos e instrucciones de acceso más frecuente para reducir el tiempo medio de acceso a memoria.

\paragraph{Respuesta 13}
\begin{itemize}
  \item E/S programada: el procesador espera activamente al dispositivo.
  \item E/S por interrupciones: el procesador realiza otras tareas y el dispositivo lo interrumpe cuando necesita atención.
  \item DMA: un controlador especializado transfiere directamente los datos entre la memoria y el dispositivo, interviniendo el procesador solo al inicio y al final.
\end{itemize}

\paragraph{Respuesta 14} Porque el procesador no participa en cada transferencia de datos; únicamente configura el controlador DMA y recibe una interrupción cuando la transferencia ha finalizado, reduciendo significativamente la carga del procesador.
